
\documentclass[../../../physics-standard_answers_all.tex]{subfiles}

\begin{document}

\setlength{\abovedisplayskip}{2mm}
\setlength{\belowdisplayskip}{1mm}

\begin{enumerate}

    \item キルヒホッフ則より，
        %
        \begin{align*}
            V_0 = \frac{q_0}{3C} \qquad \text{∴} \; q_0 = \uwave{3CV_0} \;.
        \end{align*}

    \item 電荷保存則より，
        %
        \begin{align*}
            x_1 + y_1 = q_0 + 0 \qquad \text{∴} \; x_1 + y_1 = 3CV_0 \;.
        \end{align*}
        %
        キルヒホッフ則より，
        %
        \begin{align*}
            \frac{x_1}{3C} = \frac{y_1}{2C} 
            \qquad \text{∴} \; x_1 = \uwave{\frac{9}{5} CV_0} \;,
            \quad y_1 = \uwave{\frac{6}{5} CV_0} \;.
        \end{align*}

    \item S$_2$を開き，S$_1$を閉じて充分待つと，C$_1$の電荷は再び$q_0=3CV_0$となる
        （C$_2$の電荷は$y_1$のまま）．続いて，S$_1$を開き，S$_2$を閉じて充分待った後，
        電荷保存則より，
        %
        \begin{align*}
            x_2 + y_2 = q_0 + y_1 \qquad \text{∴} \; x_2 + y_2 = \frac{21}{5}CV_0 \;.
        \end{align*}
        %
        キルヒホッフ則より，
        %
        \begin{align*}
            \frac{x_2}{3C} = \frac{y_2}{2C} 
            \qquad \text{∴} \; x_2 = \uwave{\frac{63}{25} CV_0} \;,
            \quad y_2 = \uwave{\frac{42}{25} CV_0} \;.
        \end{align*}

    \item キルヒホッフ則より，
        %
        \begin{align*}
            V_0 = \frac{X}{3C} = \frac{Y}{2C}
            \qquad \text{∴} \; X = \uwave{3CV_0} \;,
            \quad Y = \uwave{2CV_0} \;.
        \end{align*}        

\end{enumerate}

\end{document}
